% ================================================================
%  sesion03.tex  —  Sesión 3 de 20
%  Plano por cuadrantes
%  Almería en Cifras y Formas · Matemáticas 1.º ESO
%  Compilar con:  pdflatex sesion03.tex  (dos pasadas)
% ================================================================
\documentclass[aspectratio=169, 12pt, spanish]{beamer}
\usepackage{almeria-beamer}

% ---------------------------------------------------------------
%  DATOS DE LA SESIÓN
% ---------------------------------------------------------------
\renewcommand{\SessionNumber}{03}
\renewcommand{\SessionTitle}{Plano por cuadrantes}
\renewcommand{\SessionName}{Sesión 03}
\renewcommand{\SessionObjective}{%
  Localizar puntos de interés en un mapa cuadriculado de Almería usando
  un sistema de cuadrantes (letra + número) y compararlo con el sistema
  de coordenadas cartesianas.}
\renewcommand{\BlockName}{Bloque 1: Mapa y coordenadas}

% ================================================================
\begin{document}

% ---------------------------------------------------------------
%  PORTADA
% ---------------------------------------------------------------
\begin{frame}[plain]
  \titlepage
\end{frame}

% ---------------------------------------------------------------
%  ÍNDICE DE LA SESIÓN
% ---------------------------------------------------------------
\begin{frame}{Índice de la sesión}
  \begin{columns}[t]
    \begin{column}{0.5\textwidth}
      \begin{enumerate}
        \item ¿Qué es un plano cuadriculado?
        \item Cómo leer un código de cuadrante
        \item Mapa de Almería: localizar puntos
        \item Pasos para encontrar un lugar
      \end{enumerate}
    \end{column}
    \begin{column}{0.5\textwidth}
      \begin{enumerate}
        \setcounter{enumi}{4}
        \item Comparación: cuadrantes vs.\ coordenadas vs.\ GPS
        \item Ejercicios multinivel (Bloom)
        \item Evidencia del día
        \item Error frecuente
        \item Resumen
      \end{enumerate}
    \end{column}
  \end{columns}
\end{frame}

% ---------------------------------------------------------------
\seccionframe{1. ¿Qué es un plano cuadriculado?}
% ---------------------------------------------------------------

\begin{frame}{El plano cuadriculado}
  \begin{columns}[c]
    \begin{column}{0.52\textwidth}
      \begin{definicionbox}{Plano cuadriculado}
        Un mapa dividido en \textbf{filas} (letras A, B, C\ldots)
        y \textbf{columnas} (números 1, 2, 3\ldots) que forma una
        rejilla. Cada celda recibe un código único
        \textbf{letra + número}.
      \end{definicionbox}
      \vspace{6pt}
      \begin{itemize}
        \item Fácil de usar: no necesita tecnología.
        \item Estándar en guías de calles y mapas turísticos.
        \item La letra identifica la \textbf{fila horizontal}.
        \item El número identifica la \textbf{columna vertical}.
      \end{itemize}
      \vspace{4pt}
      \begin{notabox}[Regla de lectura]
        \small
        \textbf{Primero la letra} (fila),
        \textbf{después el número} (columna):
        \quad B3 = fila B, columna 3.
      \end{notabox}
    \end{column}
    \begin{column}{0.44\textwidth}
      \begin{center}
        \begin{tikzpicture}[scale=0.85]
          % Cuadrícula 4x4 básica explicativa
          \foreach \c in {1,2,3,4}{
            \foreach \f in {1,2,3,4}{
              \fill[AlmFondo, draw=AlmAzulC!50, line width=0.5pt]
                (\c-1, \f-1) rectangle (\c, \f);
            }
          }
          % Etiquetas de filas (letras)
          \foreach \f/\ltr in {4/A, 3/B, 2/C, 1/D}{
            \node[font=\small\bfseries, color=AlmAzul]
              at (-0.35, \f-0.5) {\ltr};
          }
          % Etiquetas de columnas (números)
          \foreach \c/\num in {1/1, 2/2, 3/3, 4/4}{
            \node[font=\small\bfseries, color=AlmAzul]
              at (\c-0.5, 4.35) {\num};
          }
          % Celda destacada B3
          \fill[AlmAmbar!45, draw=AlmAmbar, line width=1.5pt]
            (2, 2) rectangle (3, 3);
          \node[font=\small\bfseries, color=AlmAzul]
            at (2.5, 2.5) {B3};
          % Flechas indicadoras
          \draw[AlmTerra, ->, line width=1pt]
            (-0.35, 2.5) -- (0.05, 2.5);
          \draw[AlmAzulM, ->, line width=1pt]
            (2.5, 4.35) -- (2.5, 3.05);
        \end{tikzpicture}
      \end{center}
    \end{column}
  \end{columns}
\end{frame}

% ---------------------------------------------------------------
\seccionframe{2. Cómo leer un código de cuadrante}
% ---------------------------------------------------------------

\begin{frame}{Lectura de un código: paso a paso}
  \begin{columns}[c]
    \begin{column}{0.50\textwidth}
      \begin{pasobox}[1]
        \small
        \textbf{Lee la letra.}\\
        Identifica la \textbf{fila horizontal}
        en el margen izquierdo del mapa.
      \end{pasobox}
      \vspace{6pt}
      \begin{pasobox}[2]
        \small
        \textbf{Lee el número.}\\
        Identifica la \textbf{columna vertical}
        en el margen superior del mapa.
      \end{pasobox}
      \vspace{6pt}
      \begin{pasobox}[3]
        \small
        \textbf{Marca la intersección.}\\
        El lugar está \emph{dentro de la celda}
        donde se cruzan esa fila y esa columna.
      \end{pasobox}
    \end{column}
    \begin{column}{0.46\textwidth}
      \begin{ejemplobox}
        \small
        \textbf{Código: C4}\\[4pt]
        \begin{enumerate}
          \item Busco la fila \textbf{C}
                (tercera fila desde arriba).
          \item Busco la columna \textbf{4}
                (cuarta columna desde la izquierda).
          \item Marco donde se cruzan: ahí está el lugar.
        \end{enumerate}
      \end{ejemplobox}
      \vspace{6pt}
      \begin{formulabox}[AlmAzulM]
        \small
        Código \;=\; \textbf{letra} + \textbf{número}\\[4pt]
        \large B3 \;=\; fila\,B\;,\;columna\,3
      \end{formulabox}
    \end{column}
  \end{columns}
\end{frame}

% ---------------------------------------------------------------
\seccionframe{3. Mapa de Almería: localizar puntos}
% ---------------------------------------------------------------

\begin{frame}{Mapa cuadriculado de Almería}
  \framesubtitle{Cuadrícula 4 × 5 — puntos de interés reales}
  \begin{center}
    \begin{tikzpicture}[scale=1.05]
      % ------- Fondo del mapa -------
      \fill[AlmAzulC!12, rounded corners=4pt]
        (-0.5,-0.5) rectangle (5.5,4.5);
      % ------- Cuadrícula 4 filas × 5 columnas -------
      \foreach \c in {0,...,4}{
        \foreach \f in {0,...,3}{
          \fill[AlmBlanco!70, draw=AlmAzulC!60, line width=0.5pt]
            (\c,\f) rectangle (\c+1,\f+1);
        }
      }
      % ------- Etiquetas filas (A=top=fila 3, D=bottom=fila 0) -------
      \foreach \f/\ltr in {3/A, 2/B, 1/C, 0/D}{
        \node[font=\small\bfseries, color=AlmAzul]
          at (-0.32, \f+0.5) {\ltr};
      }
      % ------- Etiquetas columnas -------
      \foreach \c in {1,...,5}{
        \node[font=\small\bfseries, color=AlmAzul]
          at (\c-0.5, 4.28) {\c};
      }
      % ------- Puntos de interés -------
      % Plaza Catedral: B2 -> col 2, fila B(=2) -> (1.5, 2.5)
      \node[punto] (cat) at (1.5, 2.5) {};
      \node[font=\tiny\bfseries, color=AlmAzulM, above right=-1pt and 1pt of cat]
        {Catedral};
      \node[font=\tiny, color=AlmGris, below right=-1pt and 1pt of cat]
        {B2};
      % Alcazaba: A2 -> col 2, fila A(=3) -> (1.5, 3.5)
      \node[puntot] (alc) at (1.5, 3.5) {};
      \node[font=\tiny\bfseries, color=AlmTerra, above right=-1pt and 1pt of alc]
        {Alcazaba};
      \node[font=\tiny, color=AlmGris, below right=-1pt and 1pt of alc]
        {A2};
      % Mercado Central: B3 -> col 3, fila B(=2) -> (2.5, 2.5)
      \node[puntov] (mer) at (2.5, 2.5) {};
      \node[font=\tiny\bfseries, color=AlmVerde, above right=-1pt and 1pt of mer]
        {Mercado};
      \node[font=\tiny, color=AlmGris, below right=-1pt and 1pt of mer]
        {B3};
      % Puerto: C4 -> col 4, fila C(=1) -> (3.5, 1.5)
      \node[punto, fill=BAplicar] (puer) at (3.5, 1.5) {};
      \node[font=\tiny\bfseries, color=BAplicar!70!black,
            above right=-1pt and 1pt of puer]
        {Puerto};
      \node[font=\tiny, color=AlmGris, below right=-1pt and 1pt of puer]
        {C4};
      % Parque Nicolás Salmerón: C3 -> col 3, fila C(=1) -> (2.5, 1.5)
      \node[puntov] (par) at (2.5, 1.5) {};
      \node[font=\tiny\bfseries, color=AlmVerde,
            above left=-1pt and 1pt of par]
        {P. Salmerón};
      \node[font=\tiny, color=AlmGris, below left=-1pt and 1pt of par]
        {C3};
      % ------- Leyenda -------
      \node[font=\tiny, color=AlmAzul, align=left,
            draw=AlmAzulC!50, fill=AlmBlanco!80,
            rounded corners=2pt, inner sep=3pt]
        at (4.35, 3.5)
        {\textbf{Leyenda}\\
         \textcolor{AlmAzulM}{$\bullet$} Catedral\\
         \textcolor{AlmTerra}{$\bullet$} Alcazaba\\
         \textcolor{AlmVerde}{$\bullet$} Mercado\\
         \textcolor{BAplicar}{$\bullet$} Puerto\\
         \textcolor{AlmVerde}{$\bullet$} P.\,Salmerón};
    \end{tikzpicture}
  \end{center}
\end{frame}

% ---------------------------------------------------------------
\seccionframe{4. Pasos para encontrar un lugar}
% ---------------------------------------------------------------

\begin{frame}{Localización paso a paso en el mapa}
  \begin{columns}[c]
    \begin{column}{0.50\textwidth}
      \begin{teoriabox}
        \small
        Para localizar un lugar en un plano cuadriculado:
        \begin{enumerate}
          \item Busca la \textbf{letra} en el margen izquierdo
                y traza mentalmente una línea horizontal.
          \item Busca el \textbf{número} en el margen superior
                y traza una línea vertical.
          \item La celda de la \textbf{intersección}
                contiene el lugar buscado.
          \item Coloca un \textbf{punto} en el centro
                de esa celda.
        \end{enumerate}
      \end{teoriabox}
    \end{column}
    \begin{column}{0.46\textwidth}
      \begin{ejemplobox}
        \small
        \textbf{Localizar: Puerto (C4)}\\[4pt]
        \begin{enumerate}
          \item Fila \textbf{C} $\rightarrow$ tercera fila.
          \item Columna \textbf{4} $\rightarrow$ cuarta columna.
          \item Intersección $\rightarrow$ celda C4.
          \item Marca el punto \textcolor{BAplicar}{\faMapMarkerAlt}.
        \end{enumerate}
      \end{ejemplobox}
      \vspace{6pt}
      \begin{notabox}[El código siempre es único]
        \footnotesize
        Dos lugares distintos \textbf{no pueden}
        tener el mismo código.\\
        Dos códigos distintos \textbf{nunca} señalan
        al mismo lugar.
      \end{notabox}
    \end{column}
  \end{columns}
\end{frame}

% ---------------------------------------------------------------
\seccionframe{5. Comparación de sistemas de localización}
% ---------------------------------------------------------------

\begin{frame}{Cuadrantes · Coordenadas cartesianas · GPS}
  \begin{center}
    \small
    \begin{tabular}{@{}p{3.0cm}p{3.4cm}p{3.4cm}p{3.4cm}@{}}
      \toprule
      \textbf{Aspecto}
        & \textbf{Cuadrantes}\newline(letra + número)
        & \textbf{Coordenadas}\newline cartesianas $(x, y)$
        & \textbf{GPS}\newline (latitud, longitud)\\
      \midrule
      \textbf{Precisión}
        & Baja: sitúa en una celda ($\sim$km\textsuperscript{2})
        & Media: punto exacto en el mapa
        & Muy alta: $< 5$ m\\[4pt]
      \textbf{Facilidad}
        & Muy fácil;\newline sin cálculos
        & Media;\newline requiere entender ejes
        & Muy fácil con dispositivo\\[4pt]
      \textbf{Tecnología}
        & Ninguna (papel y lápiz)
        & Ninguna (cuadrícula)
        & Satélite / smartphone\\[4pt]
      \textbf{Uso típico}
        & Guías turísticas,\newline callejeros
        & Geometría, planos de arquitectura
        & Navegación, \newline cartografía digital\\[4pt]
      \textbf{Ejemplo}\newline \textbf{Almería}
        & Catedral: B2
        & Catedral: $(1{,}5;\; 2{,}5)$
        & 36,8368°\,N\newline 2,4695°\,O\\
      \bottomrule
    \end{tabular}
  \end{center}
  \vspace{4pt}
  \begin{notabox}[Conclusión]
    \footnotesize
    Cada sistema tiene su contexto ideal. Para esta SA usaremos
    \textbf{cuadrantes} y \textbf{coordenadas cartesianas}.
  \end{notabox}
\end{frame}

% ---------------------------------------------------------------
\seccionframe{6. Ejercicios multinivel — Bloom}
% ---------------------------------------------------------------

\begin{frame}{Ejercicio RECORDAR}
  \bloomtag{RECORDAR}\quad\dificultad{1}
  \vspace{4pt}
  \begin{recordarbox}
    \small
    En tu cuaderno, \textbf{dibuja} una cuadrícula 4 × 4:
    \begin{itemize}
      \item Rotula las filas con las letras A, B, C, D
            (de arriba abajo) en el margen izquierdo.
      \item Rotula las columnas con los números 1, 2, 3, 4
            (de izquierda a derecha) en el margen superior.
    \end{itemize}
    \vspace{4pt}
    Responde:
    \begin{enumerate}
      \item ¿Cuántas celdas tiene la cuadrícula en total?
      \item Colorea la celda \textbf{A3} de azul
            y la celda \textbf{D2} de rojo.
      \item ¿Qué código tiene la celda de la esquina
            superior derecha?
    \end{enumerate}
  \end{recordarbox}
\end{frame}

\begin{frame}{Ejercicio APLICAR}
  \bloomtag{APLICAR}\quad\dificultad{3}
  \vspace{4pt}
  \begin{aplicarbox}
    \footnotesize
    Usa el mapa de Almería de la diapositiva anterior.\\[4pt]
    \textbf{Parte A — Dibujar desde códigos:}
    \begin{enumerate}
      \item Localiza y marca en el mapa los siguientes lugares:\\
        Rambla de Almería: B2 \quad
        Parque Nicolás Salmerón: C3\\
        Teatro Cervantes: B3 \quad
        Playa de Almería: D4
    \end{enumerate}
    \vspace{4pt}
    \textbf{Parte B — Leer códigos del mapa:}
    \begin{enumerate}
      \setcounter{enumi}{1}
      \item El mapa muestra tres puntos sin nombre.\\
        Escribe el código de cuadrante de cada uno.
    \end{enumerate}
    \vspace{4pt}
    \textbf{Parte C — Distancia aproximada:}
    \begin{enumerate}
      \setcounter{enumi}{2}
      \item Si cada celda representa 1 km, ¿cuántos km
        hay entre la Alcazaba (A2) y el Puerto (C4)?\\
        Explica cómo has calculado esa distancia.
    \end{enumerate}
  \end{aplicarbox}
\end{frame}

\begin{frame}{Ejercicio ANALIZAR}
  \bloomtag{ANALIZAR}\quad\dificultad{4}
  \vspace{4pt}
  \begin{analizarbox}
    \small
    Compara el sistema de \textbf{cuadrantes} con el de
    \textbf{coordenadas cartesianas}.\\[6pt]
    Escribe un argumento de \textbf{5--6 líneas} que responda:
    \begin{itemize}
      \item ¿En qué situaciones es más útil el sistema de cuadrantes?
      \item ¿Cuándo preferirías coordenadas cartesianas?
      \item ¿Cuál de los dos usarías para diseñar un mapa
            turístico impreso de Almería? Justifica tu elección.
    \end{itemize}
    \vspace{4pt}
    Usa \textbf{ejemplos concretos} de Almería en tu argumento.
  \end{analizarbox}
\end{frame}

% ---------------------------------------------------------------
%  EVIDENCIA
% ---------------------------------------------------------------
\begin{frame}{Evidencia del día}
  \begin{evidenciabox}
    Entrega tu \textbf{Mapa base de Almería completado}:
    \begin{itemize}
      \item Cuadrícula de al menos 4 × 4 con líneas rectas
            y perpendiculares.
      \item Ejes correctamente nombrados
            (letras a la izquierda, números arriba).
      \item Al menos \textbf{5 puntos de interés} de Almería
            correctamente situados con su código.
      \item \textbf{Leyenda} con el nombre de cada punto.
      \item Una \textbf{frase de reflexión}:
            «El sistema de cuadrantes me sirve para\ldots»
    \end{itemize}
  \end{evidenciabox}
\end{frame}

% ---------------------------------------------------------------
%  ERROR FRECUENTE
% ---------------------------------------------------------------
\begin{frame}{¡Cuidado con este error!}
  \begin{errorbox}
    \begin{columns}[c]
      \begin{column}{0.50\textwidth}
        \incorrecto\ \textbf{Error frecuente}\\[6pt]
        Leer el código en el \textbf{orden incorrecto}:
        número antes que letra.\\[8pt]
        \begin{center}
          \large
          \textbf{\sout{3B}}~$\neq$~\textbf{B3}
        \end{center}
        \vspace{4pt}
        \footnotesize
        Si dices «3B» estás buscando la columna 3
        y la fila B, lo cual puede interpretarse
        de forma errónea y llevar a una celda
        \textbf{completamente distinta}.
      \end{column}
      \begin{column}{0.46\textwidth}
        \correcto\ \textbf{Correcto}\\[6pt]
        \textbf{Siempre primero la letra, luego el número.}\\[8pt]
        \begin{itemize}
          \item B3 $\rightarrow$ fila B, columna 3 \correcto
          \item 3B $\rightarrow$ orden incorrecto \incorrecto
        \end{itemize}
        \vspace{6pt}
        \footnotesize
        \textbf{Truco:} piensa en el alfabeto antes que en los números.
      \end{column}
    \end{columns}
  \end{errorbox}
  \vspace{6pt}
  \begin{notabox}[Analogía]
    \small
    En una hoja de cálculo, la celda B3 está
    en la columna B y la fila 3 — el mismo orden.
  \end{notabox}
\end{frame}

% ---------------------------------------------------------------
%  RESUMEN
% ---------------------------------------------------------------
\begin{frame}{Resumen de la sesión 03}
  \begin{resumenbox}
    \begin{itemize}
      \item Un \textbf{plano cuadriculado} divide el mapa en filas
            (letras) y columnas (números).
      \item El código de un lugar siempre es
            \textbf{letra primero}, número segundo: B3, C4\ldots
      \item Para localizar: \textbf{1) fila, 2) columna,
            3) intersección}.
      \item Este sistema es sencillo pero tiene
            \textbf{poca precisión} comparado con GPS.
      \item Las \textbf{coordenadas cartesianas} dan más precisión
            y las estudiaremos en la próxima sesión.
    \end{itemize}
  \end{resumenbox}
  \vspace{6pt}
  \begin{center}
    {\small\color{AlmAzulM}
    \faArrowCircleRight\enskip
    \textbf{Próxima sesión:}
    Coordenadas cartesianas — leer y escribir pares $(x, y)$.}
  \end{center}
\end{frame}

\end{document}
